
\documentclass[11pt,a4paper]{article}

\usepackage[french]{babel}
\usepackage[utf8]{inputenc}
\usepackage[T1]{fontenc}
\usepackage[a4paper,margin=2.4cm]{geometry}
\usepackage{amsmath,amssymb,amsthm,mathtools,bm}
\usepackage{enumitem}
\usepackage[hidelinks]{hyperref}
\usepackage{microtype}

\newtheorem{definition}{Définition}[section]
\newtheorem{assumption}{Hypothèse}[section]
\newtheorem{proposition}{Proposition}[section]
\newtheorem{theorem}{Théorème}[section]
\newtheorem{corollary}{Corollaire}[section]
\newtheorem{lemma}{Lemme}[section]
\newtheorem{remark}{Remarque}[section]

\DeclareMathOperator{\argmin}{argmin}
\DeclareMathOperator{\argmax}{argmax}
\DeclareMathOperator{\dist}{dist}
\DeclareMathOperator{\prox}{prox}
\DeclareMathOperator{\Tr}{tr}
\DeclareMathOperator{\Id}{Id}
\DeclareMathOperator{\op}{op}
\DeclareMathOperator{\conv}{conv}

\newcommand{\R}{\mathbb R}
\newcommand{\Hh}{\mathcal H}
\newcommand{\Mcal}{\mathcal M}
\newcommand{\Ccal}{\mathcal C}
\newcommand{\Pcal}{\mathcal P}
\newcommand{\Jcal}{\mathcal J}
\newcommand{\Ical}{\mathcal I}

\title{Extension théorique du routage contextuel dans le noyau PCH}
\author{}
\date{}

\begin{document}
\maketitle

\begin{abstract}
On prolonge la théorie PCH déjà établie en isolant explicitement la couche de routage.
Le noyau PCH n'est pas modifié : pour chaque contexte admissible, on conserve l'opérateur exact et sa normalisation.
La nouvelle contribution est une fonctionnelle de sélection de contexte fondée sur l'évidence variationnelle normalisée, sur un certificat de gain, puis sur un raffinement géométrique par projecteur limite.
Cette construction sépare la purification locale exacte et le choix de contexte, ce qui retire au routage son rôle de goulot d'étranglement sans altérer la dynamique PCH elle-même.
\end{abstract}

\section{Préliminaires et jauge normalisée}

On se place dans la même jauge que celle du Cas 3 normalisé.
Pour un index de famille $\alpha$, on note
\[
\mu_\alpha := 1+2\lambda_{\mathrm{ridge},\alpha},\qquad
\gamma_\alpha := \frac{\lambda_{\mathrm{ctx},\alpha}}{\mu_\alpha},\qquad
\tau_\alpha := \eta_\alpha \mu_\alpha.
\]
Sous la normalisation retenue dans la théorie précédente, les quantités $\gamma_\alpha$ et $\tau_\alpha$ convergent vers des constantes déterministes, et le projecteur différentiel local se concentre vers un projecteur orthogonal limite $P_0$.

Pour chaque contexte $c\in\Ccal$, on dispose d'un ensemble admissible convexe fermé
\[
\Mcal_{\alpha,c}\subset \Hh^{d_\alpha}
\]
construit à partir des embeddings de calibration.
L'énergie locale normalisée est
\[
\widehat J_{\alpha,c}(u;w)
=
\frac{1}{\mu_\alpha}
\Bigl(
\frac12\|u-w\|^2
+
\lambda_{\mathrm{ctx},\alpha}\,\dist(u,\Mcal_{\alpha,c})^2
+
\lambda_{\mathrm{ridge},\alpha}\,\|u\|^2
\Bigr).
\]
L'opérateur PCH exact du contexte $c$ est
\[
u^\star_{\alpha,c}(w)
=
\Pcal^{\mathrm{PCH}}_{\alpha,c}(w)
=
\argmin_{u\in \Hh^{d_\alpha}} \widehat J_{\alpha,c}(u;w).
\]
Le noyau PCH reste donc strictement inchangé : seule la couche de décision entre contextes sera définie plus bas.

\section{Pourquoi le routage doit être formulé comme un problème variationnel}

Le choix de contexte ne doit pas reposer sur un classifieur plat $x\mapsto c$.
Il doit comparer, pour un même \(w\), les minima exacts des problèmes PCH associés à chaque contexte.
La quantité brute pertinente pour un contexte \(c\) est le gain exact
\begin{equation}\label{eq:gain-def}
\Delta_{\alpha,c}(w)
=
\widehat J_{\alpha,c}(w;w)
-
\widehat J_{\alpha,c}\bigl(u^\star_{\alpha,c}(w);w\bigr).
\end{equation}
Elle mesure la baisse d'énergie obtenue lorsque le point d'entrée est purifié par le contexte \(c\).

\begin{proposition}[Invariance par recentrage additif]
\label{prop:additive-invariance}
Soit \(b_{\alpha,c}:\Hh^{d_\alpha}\to \R\) une fonction ne dépendant pas de \(u\), et posons
\[
\widetilde J_{\alpha,c}(u;w)=\widehat J_{\alpha,c}(u;w)+b_{\alpha,c}(w).
\]
Alors le gain associé satisfait
\[
\widetilde\Delta_{\alpha,c}(w)=\Delta_{\alpha,c}(w).
\]
\end{proposition}

\begin{proof}
Les termes ajoutés s'annulent dans la différence
\[
\widetilde J_{\alpha,c}(w;w)-\widetilde J_{\alpha,c}(u^\star_{\alpha,c}(w);w)
=
\widehat J_{\alpha,c}(w;w)-\widehat J_{\alpha,c}(u^\star_{\alpha,c}(w);w).
\]
Donc le gain est invariant sous tout recentrage additif contexte-dépendant.
\end{proof}

\begin{remark}
Cette invariance est précisément ce qui manque aux comparaisons par énergie finale brute lorsqu'elles sont appliquées de façon naïve entre contextes.
Le bon objet n'est pas la valeur absolue du minimum, mais la quantité d'information effectivement extraite par la purification.
\end{remark}

\section{Score canonique de routage}

\begin{definition}[Score de routage par gain normalisé]
Pour un contexte \(c\), on définit
\begin{equation}\label{eq:routing-score}
\mathcal R_{\alpha,c}(w)
=
\frac{\Delta_{\alpha,c}(w)}{\widehat J_{\alpha,c}(w;w)+\varepsilon},
\qquad \varepsilon>0.
\end{equation}
\end{definition}

\begin{proposition}[Bornes]
Pour tout \(w\) et tout \(c\),
\[
0\le \mathcal R_{\alpha,c}(w)\le 1.
\]
De plus, \(\mathcal R_{\alpha,c}(w)=1\) si et seulement si \(w\) est déjà un minimiseur de \(\widehat J_{\alpha,c}(\cdot;w)\).
\end{proposition}

\begin{proof}
Par définition, \(\Delta_{\alpha,c}(w)\ge 0\) puisque \(u^\star_{\alpha,c}(w)\) minimise \(\widehat J_{\alpha,c}(\cdot;w)\).
De plus,
\[
0\le \Delta_{\alpha,c}(w)\le \widehat J_{\alpha,c}(w;w),
\]
puisque \(\widehat J_{\alpha,c}(u^\star_{\alpha,c}(w);w)\ge 0\).
La division par un terme strictement positif préserve donc l'encadrement.
L'égalité \(\mathcal R_{\alpha,c}(w)=1\) équivaut à \(\widehat J_{\alpha,c}(u^\star_{\alpha,c}(w);w)=0\), donc à l'optimalité immédiate de \(w\).
\end{proof}

\begin{definition}[Routeur exact par evidence]
On définit le contexte retenu par
\begin{equation}\label{eq:hard-router}
c^\star_\alpha(w)\in
\argmax_{c\in\Ccal}\mathcal R_{\alpha,c}(w).
\end{equation}
Le point purifié final est alors
\[
\widetilde u_\alpha(w)=u^\star_{\alpha,c^\star_\alpha(w)}(w).
\]
\end{definition}

\begin{theorem}[Routage exact sans modification du noyau]
La règle \eqref{eq:hard-router} ne modifie en rien le noyau PCH :
pour chaque contexte \(c\), on résout le même problème exact qu'avant.
La seule nouveauté est l'agrégation externe qui choisit le contexte maximisant le gain normalisé.
\end{theorem}

\begin{proof}
Pour tout \(c\), le calcul de \(u^\star_{\alpha,c}(w)\) est identique à celui de la purification contextuelle déjà définie.
La couche de routage n'ajoute qu'une comparaison scalaire entre les gains \(\mathcal R_{\alpha,c}(w)\).
Aucun terme du noyau PCH n'est modifié.
\end{proof}

\begin{theorem}[Stabilité de la décision sous marge]
\label{thm:margin-stability}
Supposons qu'il existe un contexte \(c^\star\) et une marge \(\delta>0\) tels que
\[
\mathcal R_{\alpha,c^\star}(w)\ge \mathcal R_{\alpha,c}(w)+2\delta
\qquad\text{pour tout }c\neq c^\star.
\]
Si des perturbations \(\eta_c(w)\) satisfont
\[
|\eta_c(w)|\le \delta\qquad\text{pour tout }c,
\]
alors le contexte sélectionné par les scores perturbés
\(\mathcal R_{\alpha,c}(w)+\eta_c(w)\) reste \(c^\star\).
\end{theorem}

\begin{proof}
Pour tout \(c\neq c^\star\),
\[
\mathcal R_{\alpha,c^\star}(w)+\eta_{c^\star}(w)
\ge
\mathcal R_{\alpha,c^\star}(w)-\delta
\ge
\mathcal R_{\alpha,c}(w)+\delta
\ge
\mathcal R_{\alpha,c}(w)+\eta_c(w).
\]
Le maximiseur reste donc inchangé.
\end{proof}

\section{Routage probabiliste et limite dure}

Un routeur probabiliste naturel est donné par la loi de Gibbs
\begin{equation}\label{eq:soft-router}
\pi_{\beta,\alpha}(c\mid w)
=
\frac{\exp\!\bigl(\beta\,\mathcal R_{\alpha,c}(w)\bigr)}
{\sum_{j\in\Ccal}\exp\!\bigl(\beta\,\mathcal R_{\alpha,j}(w)\bigr)},
\qquad \beta>0.
\end{equation}

\begin{theorem}[Concentration vers le routeur dur]
\label{thm:soft-to-hard}
Si \(c^\star\) est l'unique maximiseur de \(\mathcal R_{\alpha,c}(w)\), alors
\[
\pi_{\beta,\alpha}(c^\star\mid w)\xrightarrow[\beta\to\infty]{}1
\qquad\text{et}\qquad
\pi_{\beta,\alpha}(c\mid w)\xrightarrow[\beta\to\infty]{}0
\quad(c\neq c^\star).
\]
\end{theorem}

\begin{proof}
Écrivons \(\Delta_c=\mathcal R_{\alpha,c^\star}(w)-\mathcal R_{\alpha,c}(w)\), avec \(\Delta_c>0\) pour \(c\neq c^\star\).
Alors
\[
\pi_{\beta,\alpha}(c^\star\mid w)
=
\frac{1}{1+\sum_{c\neq c^\star}e^{-\beta\Delta_c}}
\to 1.
\]
Le reste suit immédiatement.
\end{proof}

\begin{corollary}
Le routeur dur \eqref{eq:hard-router} est la limite froide du routeur probabiliste \eqref{eq:soft-router}.
Ainsi, la décision exacte peut être lue comme une sélection de type \(\beta\to\infty\).
\end{corollary}

\section{Certificat de routage par pré-roulage court}

La théorie permet un routeur exact, mais son coût peut être élevé.
On introduit donc un certificat approché fondé sur un pré-roulage de \(K\) itérations par contexte.

Pour chaque contexte \(c\), on définit
\[
u^{(0)}_{\alpha,c}(w)=w,
\qquad
u^{(k+1)}_{\alpha,c}(w)
=
\mathcal T_{\eta_\alpha,\alpha}\!\bigl(u^{(k)}_{\alpha,c}(w);w,c\bigr),
\qquad k\ge 0,
\]
où \(\mathcal T_{\eta_\alpha,\alpha}\) est l'opérateur proximal exact du contexte \(c\).
Le certificat court est alors
\[
\Delta^{(K)}_{\alpha,c}(w)
=
\widehat J_{\alpha,c}(w;w)
-
\widehat J_{\alpha,c}\bigl(u^{(K)}_{\alpha,c}(w);w\bigr),
\]
et son score normalisé
\[
\mathcal R^{(K)}_{\alpha,c}(w)
=
\frac{\Delta^{(K)}_{\alpha,c}(w)}{\widehat J_{\alpha,c}(w;w)+\varepsilon}.
\]

\begin{proposition}[Monotonie du certificat court]
Sous forte convexité locale, la suite \(k\mapsto \widehat J_{\alpha,c}(u^{(k)}_{\alpha,c}(w);w)\) est décroissante.
Par conséquent,
\[
0\le \Delta^{(1)}_{\alpha,c}(w)\le \cdots \le \Delta^{(K)}_{\alpha,c}(w)\le \Delta_{\alpha,c}(w),
\]
et les scores correspondants vérifient la même monotonie.
\end{proposition}

\begin{proof}
Chaque itération est définie comme un minimiseur proximal.
L'inégalité de décroissance du fonctionnel proximal donne la monotonie de l'énergie.
En soustrayant à la valeur initiale, on obtient la monotonie du gain.
\end{proof}

\begin{theorem}[Contrôle géométrique de l'erreur de routage court]
\label{thm:short-rollout}
Supposons que, pour chaque contexte \(c\), l'opérateur proximal local soit contractant avec facteur \(\rho_c\in(0,1)\).
Supposons en outre qu'il existe une constante \(L_c>0\) telle que le score \(\mathcal R_{\alpha,c}\) soit \(L_c\)-Lipschitz par rapport à l'état latent sur la trajectoire de pré-roulage.
Alors il existe une constante \(C_c>0\) telle que
\[
\bigl|\mathcal R^{(K)}_{\alpha,c}(w)-\mathcal R_{\alpha,c}(w)\bigr|
\le
C_c\,\rho_c^K.
\]
\end{theorem}

\begin{proof}
La contraction donne
\[
\|u^{(K)}_{\alpha,c}(w)-u^\star_{\alpha,c}(w)\|
\le
\rho_c^K\,\|w-u^\star_{\alpha,c}(w)\|.
\]
Par Lipschitzianité du score, on en déduit
\[
\bigl|\mathcal R^{(K)}_{\alpha,c}(w)-\mathcal R_{\alpha,c}(w)\bigr|
\le
L_c\|u^{(K)}_{\alpha,c}(w)-u^\star_{\alpha,c}(w)\|.
\]
En posant \(C_c=L_c\|w-u^\star_{\alpha,c}(w)\|\), on obtient l'estimation annoncée.
\end{proof}

\begin{corollary}[Exactitude finie sous marge]
\label{cor:finite-exactness}
Si le contexte optimal \(c^\star\) vérifie
\[
\mathcal R_{\alpha,c^\star}(w)\ge \mathcal R_{\alpha,c}(w)+2\delta
\qquad(c\neq c^\star)
\]
et si \(K\) est tel que
\[
\max_{c\in\Ccal}\bigl|\mathcal R^{(K)}_{\alpha,c}(w)-\mathcal R_{\alpha,c}(w)\bigr|<\delta,
\]
alors le routage court renvoie le même contexte que le routage exact.
\end{corollary}

\begin{proof}
Cela résulte immédiatement du théorème de stabilité par marge, en prenant comme perturbations les erreurs de troncature du pré-roulage.
\end{proof}

\section{Raffinement géométrique et tie-breaker universel}

Le score de gain normalisé résout le choix principal.
Lorsque plusieurs contextes sont quasi ex æquo, il faut un départage géométrique compatible avec la normalisation universelle.

On note
\[
P_{\alpha,c}
=
D\Pi_{\Mcal_{\alpha,c}}\bigl(u^\star_{\alpha,c}(w)\bigr)
\]
le différentiel local de projection au point fixe du contexte \(c\).
La théorie normalisée fournit un projecteur limite \(P_0\).

\begin{definition}[Divergence géométrique normalisée]
On pose
\[
\Xi_{\alpha,c}(w)
=
-\bigl\|P_{\alpha,c}-P_0\bigr\|_{\op}.
\]
\end{definition}

\begin{definition}[Routage lexicographique]
Le routage final choisit le contexte maximisant le couple
\[
\bigl(\mathcal R_{\alpha,c}(w),\ \Xi_{\alpha,c}(w)\bigr)
\]
au sens lexicographique :
on maximise d'abord \(\mathcal R_{\alpha,c}(w)\), puis, en cas d'égalité ou de quasi-égalité numérique, on maximise \(\Xi_{\alpha,c}(w)\), c'est-à-dire qu'on minimise la distance au projecteur universel \(P_0\).
\end{definition}

\begin{proposition}[Cohérence du tie-breaker]
Si deux contextes sont à gain égal, le tie-breaker géométrique sélectionne celui dont la géométrie locale est la plus proche du régime universel normalisé.
Ce départage est compatible avec la théorie du Cas 3 normalisé, car il ne modifie ni la fonctionnelle PCH ni la jauge \((\mu_\alpha,\gamma_\alpha,\tau_\alpha)\).
\end{proposition}

\begin{proof}
Le critère géométrique n'intervient qu'après égalité du score principal.
Il ne touche donc pas à la définition du noyau PCH.
De plus, il compare directement les différentielles de projection au projecteur limite \(P_0\), ce qui est précisément la quantité contrôlée dans la théorie normalisée.
\end{proof}

\begin{theorem}[Sélection finale canonique]
La règle combinée
\[
c^\star_\alpha(w)
=
\argmax_{c\in\Ccal}
\Bigl(\mathcal R_{\alpha,c}(w),\ \Xi_{\alpha,c}(w)\Bigr)
\]
est un routeur externe canonique pour PCH.
Elle laisse le noyau intact, compare les contextes sur une quantité d'évidence invariante par recentrage additif, et ne mobilise le projecteur limite universel qu'en second rang.
\end{theorem}

\section{Lecture opérationnelle}

La couche de sélection peut être résumée ainsi :
\[
\boxed{
\begin{aligned}
u^\star_{\alpha,c}(w)&=\Pcal^{\mathrm{PCH}}_{\alpha,c}(w),\\[0.2em]
\mathcal R_{\alpha,c}(w)&=
\frac{
\widehat J_{\alpha,c}(w;w)
-
\widehat J_{\alpha,c}(u^\star_{\alpha,c}(w);w)
}{
\widehat J_{\alpha,c}(w;w)+\varepsilon
},\\[0.2em]
c^\star_\alpha(w)&=\argmax_{c\in\Ccal}
\bigl(\mathcal R_{\alpha,c}(w),\Xi_{\alpha,c}(w)\bigr),\\[0.2em]
\widetilde u_\alpha(w)&=u^\star_{\alpha,c^\star_\alpha(w)}(w).
\end{aligned}
}
\]

Cette écriture montre clairement que :
\begin{itemize}[leftmargin=2em]
\item le noyau PCH est entièrement préservé ;
\item le choix de contexte est formulé comme une sélection variationnelle normalisée ;
\item la version courte par pré-roulage fournit un certificat d'approximation ;
\item le tie-breaker géométrique utilise uniquement la théorie normalisée déjà validée.
\end{itemize}

\begin{remark}
Le rôle du routage n'est donc plus de deviner un label de contexte par apprentissage direct.
Il devient celui d'un problème de sélection de modèles certifié par le gain d'énergie, puis départagé par la géométrie universelle lorsque nécessaire.
C'est précisément ce qui retire au routage son statut de goulot d'étranglement sans altérer le noyau PCH.
\end{remark}

\end{document}
